\thispagestyle{empty}
\chapter*{กิตติกรรมประกาศ}
\addcontentsline{toc}{chapter}{กิตติกรรมประกาศ}

\noindent ตำราวิชาการเรื่อง \textbf{OpenXR และการประมวลผลเชิงพื้นที่เพื่อสะเต็มศึกษา} เล่มนี้สำเร็จลุล่วงได้ด้วยความกรุณาสนับสนุนและความช่วยเหลืออย่างจริงใจจากหลายท่านและหลายองค์กร ผู้เขียนขอขอบพระคุณอย่างสูงมา ณ โอกาสนี้

\noindent ขอขอบพระคุณ \textbf{มหาวิทยาลัยราชภัฏรำไพพรรณี} ที่ให้ทุนสนับสนุนการจัดทำตำราวิชาการและให้เวลาในการวิจัยและพัฒนาเนื้อหา รวมถึงคณะวิทยาศาสตร์และเทคโนโลยีที่ให้การสนับสนุนทางทรัพยากรและห้องปฏิบัติการ

\noindent ขอขอบพระคุณ \textbf{คณาจารย์และนักวิชาการ} ในสาขาวิชาฟิสิกส์ วิทยาการคอมพิวเตอร์ และวิศวกรรมซอฟต์แวร์ ทั้งในและต่างประเทศ ที่ให้คำปรึกษาและข้อเสนอแนะอันทรงคุณค่า โดยเฉพาะด้านมาตรฐาน Khronos OpenXR และการจำลองทางฟิสิกส์เชิงคำนวณ

\noindent ขอขอบพระคุณ \textbf{นักศึกษาระดับปริญญาตรีและบัณฑิตศึกษา} ที่ร่วมทดสอบชุดปฏิบัติการ Python และ OpenXR เสมือนจริง และให้ข้อเสนอแนะอันเป็นประโยชน์ยิ่งในการปรับปรุงเนื้อหาให้มีความถูกต้องและมีความเป็นประโยชน์ต่อการเรียนรู้อย่างแท้จริง

\noindent ขอขอบพระคุณ \textbf{Khronos Group} ผู้พัฒนาและเผยแพร่เอกสารมาตรฐาน OpenXR Specification รวมถึงชุมชนนักพัฒนาโอเพนซอร์ส ได้แก่ Monado, OpenXR SDK, และ PyOpenXR ที่ทำให้การวิจัยและพัฒนาตำรานี้เป็นไปได้

\noindent สุดท้ายนี้ ขอขอบพระคุณ \textbf{ครอบครัว} ที่ให้กำลังใจและความอดทนตลอดระยะเวลาการจัดทำตำราเล่มนี้ด้วยความรักและความเข้าใจ ผู้เขียนหวังเป็นอย่างยิ่งว่าตำราเล่มนี้จะเป็นประโยชน์ต่อผู้เรียนและวงการวิชาการสะเต็มศึกษาของประเทศไทยต่อไป

\vspace{1.5cm}
\begin{flushright}
    {\bfseries\color{spatialBlue} ผู้ช่วยศาสตราจารย์ ดร.ชีวะ ทัศนา}\\[4pt]
    สาขาวิชาฟิสิกส์ คณะวิทยาศาสตร์และเทคโนโลยี\\
    มหาวิทยาลัยราชภัฏรำไพพรรณี\\[4pt]
    {\color{slateGrey} พุทธศักราช 2569}
\end{flushright}
\clearpage
